\documentclass[tikz,border=4pt]{standalone}
\usepackage{ctex}
\usepackage{amsmath}
\usetikzlibrary{calc}

\begin{document}
% part12/02 两骰子列表法（6×6表格，高亮和=7的对角线）
\begin{tikzpicture}[scale=0.72, font=\small]
  % 表格：行=第一枚骰子(1-6), 列=第二枚骰子(1-6)
  % 格子大小 0.72cm
  \def\sz{1}

  % 高亮和=7的6个格子（背景）
  \foreach \r/\c in {1/6, 2/5, 3/4, 4/3, 5/2, 6/1}{
    \fill[orange!30] ({(\c)*\sz}, {-(\r)*\sz}) rectangle ({(\c+1)*\sz}, {-(\r-1)*\sz});
  }

  % 网格线
  \draw[gray] (0,0) grid (7,-6);
  \draw[thick] (0,0) rectangle (7,-6);
  % 表头列（第二枚骰子）
  \foreach \c in {1,...,6}{
    \node[font=\footnotesize] at ({\c*\sz+0.5}, 0.35) {$\c$};
  }
  % 表头行（第一枚骰子）
  \foreach \r in {1,...,6}{
    \node[font=\footnotesize] at (-0.35, {-(\r-0.5)*\sz}) {$\r$};
  }
  % 表头标签
  \node[font=\small, rotate=90] at (-0.8, -3) {第一枚};
  \node[font=\small] at (3.5, 0.75) {第二枚骰子};

  % 填入和的数值
  \foreach \r in {1,...,6}{
    \foreach \c in {1,...,6}{
      \pgfmathsetmacro{\s}{int(\r+\c)}
      \pgfmathparse{int(\r+\c)==7 ? 1 : 0}
      \ifnum\pgfmathresult=1
        \node[red, font=\small\bfseries] at ({(\c)*\sz+0.5}, {-(\r-0.5)*\sz}) {$\s$};
      \else
        \node[font=\footnotesize] at ({(\c)*\sz+0.5}, {-(\r-0.5)*\sz}) {$\s$};
      \fi
    }
  }

  % 图题
  \node[font=\small] at (3.5, -6.6) {两骰子点数之和（和$=7$用橙色高亮）};
  \node[font=\small] at (3.5, -7.1) {$P(\text{和}=7)=\dfrac{6}{36}=\dfrac{1}{6}$};
\end{tikzpicture}
\end{document}
